\section*{Author Section 6: PSDness of the exact pseudo-expectation}
Henceforth $M$ exclusively refers to the $d/2$-homogeneous minor of the moment matrix $\widetilde M$ in Definition3.13.
\paragraph{Theorem 6.1.} W.p. $>1-n^{-5\log n}$,
\[M(G)\succeq n^{-d-1}\operatorname{diag}(\widetilde{\mathrm{Cl}}(G))_{\binom{[n]}{d/2}\times\binom{[n]}{d/2}}.\]
\paragraph{Corollary 6.2.} W.p. $>1-n^{-5\log n}$, $\widetilde{\mathbb E}x_\emptyset>0$.
\begin{proof}By construction(3.13),
\[\widetilde{\mathbb E}x_\emptyset=\frac{\binom{\omega-d/2}{d-d/2}}{\binom{\omega-d}d\binom d{d/2}}\sum_{S:|S|=d/2}\widetilde{\mathbb E}x_S=\frac{\binom{\omega-d/2}{d-d/2}}{\binom{\omega-d}d\binom d{d/2}}\operatorname{Tr}(M),\]
positive w.h.p. by Theorem6.1.\end{proof}
Theorem1.4 is a quick corollary: the pseudo-expectation from Definition3.13 has PSD moment matrix by Theorem6.1 and Lemma4.1; Default holds by Corollary6.2 and rescaling as discussed before Remark3.9; Clique and Size Constraints hold by Lemma3.8. The degree-$d$ lower bound follows. The remainder proves Theorem6.1 by reducing it to Main Lemma6.9 and proving that lemma.
\subsection*{6.1 A Hadamard product and Euler transform}
In $M(I,J)$ the parameter $u=|I\cap J|$ prevents termwise factorization like the non-exact case. We express $M$ in a $\Sigma\Pi$-form where leaf matrices no longer depend on $u$, and factor those later.
\subsubsection*{6.1.1 Hadamard product}
Put $u=|I\cap J|$, $a=|V(T)\cup I\cup J|$. From(3.17),
\[M(I,J;T)=\sum_{c=0}^u\left[\frac{\omega^{u-c}}{\binom{\omega-d+u}u}\mathcal M_c(u,a)\right],\tag{6.1}\]
where
\[\mathcal M_c(u,a):=\binom{a-(d-u)}c\binom{n-a}{u-c}n^{-(u-c)}\frac{(a+u-c+8\tau^2)!}{(8\tau^2)!}(\omega/n)^a.\]
Thus
\[M=\sum_{c=0}^{d/2}m_c\circ M_c,\tag{6.2}\]
with matrices, for $|I|=|J|=d/2$,
\[m_c(I,J)=\frac{\omega^{u-c}}{\binom{\omega-d+u}u},\qquad u=|I\cap J|,\tag{6.3}\]
and
\[M_c(I,J)=\begin{cases}\sum_{T:|V(T)\cup I\cup J|\le\tau}\chi_T\mathcal M_c(|I\cap J|,|V(T)\cup I\cup J|),&|I\cap J|\ge c,\\0,&\text{o.w.}\end{cases}\tag{6.4}\]
\paragraph{Remark 6.3.} $m_c$ is supported on all $(I,J)$, while $M_c(I,J)=0$ when $|I\cap J|<c$. The intuition is that $M_c$'s are ``close'' for varying $c$, while $m_c$ qualitatively decreases in $c$. This would let us concentrate on $c=0$. The next lemma proves the second half; Lemma6.9 states the first more precisely.
\paragraph{Lemma 6.4.} For $c=0,\ldots,d/2$,
\[m_c=\omega^{-c}\sum_{k=0}^{d/2}b_kJ_k,\qquad b_k/k!\in[d/(2\omega),1+2dk/\omega],\]
where $J_k$ is (5.4). In particular,
\[m_0=\omega m_1=\cdots=\omega^{d/2}m_{d/2}\succ\frac1\omega\mathrm{Id}.\tag{6.5}\]
\begin{proof}By definition,
\[m_c=\omega^{-c}\sum_{l=0}^{d/2}\frac{\omega^l}{\binom{\omega-d+l}l}D_l.\]
By basis-change(5.6), $m_c=\omega^{-c}\sum_{k=0}^{d/2}J_kb_k$ with
\[b_k=k!\left(\sum_{l=0}^k(-1)^{k-l}\underbrace{\frac\omega{\omega-(d-l)}\cdots\frac\omega{\omega-(d-1)}\frac1{(k-l)!}}_{=:f_k(l),\text{ which is }1/k!\text{ if }l=0}\right).\]
For fixed $k$, $f_k(l)$ is increasing, so
\[\sum_{l=0}^k(-1)^{k-l}f_k(l)\ge f_k(k)-f_k(k-1)>\frac{d/2}\omega(1+(d/2)/\omega)^{k-1}\ge\frac d{2\omega}.\]
For $k=d/2$, $J_{d/2}=\mathrm{Id}$ gives(6.5).\end{proof}
\subsubsection*{6.1.2 Euler transform}
Fix $c$. For $(I,J;T)$ the coefficient $\mathcal M_c(u,a)$ is the partial function
\[\mathcal M_c(u,a)=\binom{a-(d-u)}c\binom{n-a}{u-c}n^{-(u-c)}\frac{(a+u-c+8\tau^2)!}{(8\tau^2)!}(\omega/n)^a.\tag{6.6}\]
\paragraph{Definition 6.5 (Extended $\mathcal M_c$).} It is defined on $u\ge c$, $u+a\ge d+c$ and extended to $\mathbb N^2$ by
\[\binom{n-a}{u-c}=0\text{ if }u<c,\tag{6.7}\]
and the usual conventions
\[\binom{-m}k=(-1)^k\binom{m+k-1}k\quad(0<m,\ 0\le k),\tag{6.8}\]
\[\binom mk=0\quad(0\le m<k).\tag{6.9}\]
In particular $\binom m0=1$ for all $m\in\mathbb Z$, and if $0\le a-(d-u)<c$, then $\mathcal M_c(u,a)=0$.
Seek a decomposition
\[M_c=\sum_{R\in\binom{[n]}{\le d/2}}M_c^R,\tag{6.10}\]
where $M_c^R$ is supported on row/column indices containing $R$. Explicitly,
\[M_c^R(I,J)=\begin{cases}\sum_{T:|V(T)\cup I\cup J|\le\tau}(\omega/n)^aY_c(|R|,a)\chi_T,&R\subseteq I,J,\\0,&\text{o.w.}\end{cases}\tag{6.11}\]
where $a=|V(T)\cup I\cup J|$ and $Y_c$ is to be chosen. Comparing each $(I,J;T)$, (6.10) is equivalent, for fixed $c,a$, to
\[\sum_{r=0}^u\binom urY_c(r,a)(\omega/n)^a=\mathcal M_c(u,a).\tag{6.12}\]
Thus take $Y_c(u,a)(\omega/n)^a$ to be the inverse Euler transform in $u$ of the extended $\mathcal M_c$.
\paragraph{Fact 6.6.} If sequences on $\mathbb N$ satisfy $x(m)=\sum_{l=0}^m\binom ml y(l)$, then $x$ is the Euler transform of $y$, whose inverse is $y(m)=\sum_{l=0}^m(-1)^{m-l}\binom ml x(l)$.\footnote{This fact is an application of the $\zeta$ matrix and its inverse.}
\paragraph{Definition 6.7 (Coefficients in $M_c^R$).}
\[Y_c(r,a)=\begin{cases}\sum_{l=c}^r(-1)^{r-l}\binom rl\binom{a+l-d}c\binom{n-a}{l-c}n^{-(l-c)}\dfrac{(a+l-c+8\tau^2)!}{(8\tau^2)!},&r\ge c,\\0,&\text{o.w.}\end{cases}\tag{6.13}\]
\paragraph{Lemma 6.8 (Hadamard-product decomposition).}
\[M=\sum_{c=0}^{d/2}m_c\circ\left(\sum_RM_c^R\right)=\sum_{R\in\binom{[n]}{\le d/2}}\underbrace{\left(\sum_{c=0}^{|R|}m_c\circ M_c^R\right)}_{=:M^R}.\tag{6.14--6.15}\]
Each $m_c$ is Lemma6.4, and:
\begin{enumerate}\item $M_c^R=0$ if $|R|<c$.
\item If $R\not\subseteq I\cap J$, $M_c^R(I,J)=0$.
\item If $|R|\ge c$ and $R\subseteq I\cap J$, then $M_c^R(I,J)=\sum_{T:|V(T)\cup I\cup J|\le\tau}M_c^R(I,J;T)\chi_T$, with
\begin{align*}M_c^R(I,J;T)&=(\omega/n)^a\sum_{l=c}^{|R|}(-1)^{|R|-l}\binom{|R|}l\binom{a+l-d}c\\&\quad\cdot\binom{n-a}{l-c}n^{-(l-c)}\frac{(a+l-c+8\tau^2)!}{(8\tau^2)!}\\&=(\omega/n)^aY_c(|R|,a).\tag{6.16}\end{align*}
\item For $0\le c\le r\le d/2$, $0\le a\le\tau$, $|Y_c(r,a)|<\tau^{5\tau}$.\end{enumerate}
\begin{proof}(1),(2),(3) are definition. To check $M_c=\sum_RM_c^R$, fix $(I,J;T)$ with $|I|=|J|=d/2$, $a=|V(T)\cup I\cup J|\le\tau$, and $u=|I\cap J|$. Note $a-(d-u)\ge0$ and
\[\sum_RM_c^R(I,J;T)=\sum_{R\subseteq I\cap J}M_c^R(I,J;T)=(\omega/n)^a\sum_{r=0}^{|I\cap J|}\binom{|I\cap J|}rY_c(r,a).\]
By Euler transform and(6.12) this equals extended $\mathcal M_c(u,a)$. It only remains to see $\mathcal M_c(u,a)=0$ if $u<c$ or $a-(d-u)<c$, by(6.7),(6.9). For(4),
\begin{align*}|Y_c(u,a)|&=\left|\sum_{l=c}^r(-1)^{r-l}\binom rl\binom{a+l-d}c\left[\binom{n-a}{l-c}n^{-(l-c)}\right]\frac{(a+l-c+8\tau^2)!}{(8\tau^2)!}\right|\\&<r\,2^r(2\tau)^r\cdot1\cdot(9\tau^2)^{2\tau}<\tau^{5\tau},\end{align*}
using $r\le d/2\ll\tau$ in our regime.\end{proof}
\paragraph{Lemma 6.9 (Main Lemma).} In(6.15), w.p. $>1-n^{-5\log n}$, for every $R\in\binom{[n]}{\le d/2}$, with $\mathcal P^R=\{I\in\binom{[n]}{d/2}:R\subseteq I\}$,
\[M_0^R\succeq n^{-d}\operatorname{diag}(\widetilde{\mathrm{Cl}})_{\mathcal P^R\times\mathcal P^R},\tag{6.17}\]
and
\[\pm\omega^{-c}M_c^R\preceq n^{-c/6}M_0^R\quad\forall0<c\le|R|.\tag{6.18}\]
\paragraph{Corollary 6.10 (Theorem6.1).} W.p. $>1-n^{-5\log n}$, $M(G)\succeq n^{-d-1}\operatorname{diag}(\widetilde{\mathrm{Cl}}(G))$ on $\binom{[n]}{d/2}\times\binom{[n]}{d/2}$.
\begin{proof}Suppose Lemma6.9 holds. For fixed $R$, $M^R=\sum_{c=0}^{|R|}m_c\circ M_c^R$. Hadamard product with a PSD matrix preserves PSDness (Schur product theorem), so
\begin{align*}\sum_{c=1}^{|R|}m_c\circ M_c^R&\preceq\sum_{c=1}^{|R|}m_c\circ(\omega^cn^{-c/6}M_0^R)\\&=\left(\sum_{c=1}^{|R|}n^{-c/6}m_0\right)\circ M_0^R\\&\preceq n^{-1/6}m_0\circ M_0^R.\end{align*}
Similarly the sum is $\succeq-n^{-1/6}m_0\circ M_c^R$. Thus
\[M^R\succeq(1-n^{-1/6})m_0\circ M_0^R\succeq n^{-d-1}\operatorname{diag}(\widetilde{\mathrm{Cl}})_{\mathcal P^R\times\mathcal P^R}\]
by Lemmas6.4/6.9(2). Apply(6.15):
\[M=M^\emptyset+\sum_{\emptyset\ne R\in\binom{[n]}{\le d/2}}M^R\succeq M^\emptyset\succeq n^{-d-1}\operatorname{diag}(\widetilde{\mathrm{Cl}})_{\binom{[n]}{d/2}\times\binom{[n]}{d/2}}.\tag{6.19}\]\end{proof}
